% Follow-up supplement. Standalone finding IDs reserved by root: A07,A09,A10,A11.
% A08 is supplied separately by effective_galois/finite_prime_avoidance.tex.
% Bibliography additions are in POSITIVE_RANK_BIB.tex.

\section{Positive rank and a weaker geometric bridge}
\label{sec:positive-rank-followup}

\paragraph{Finding A07 (the generic positive-rank bridge fails).}
For each $f\in\{1,2\}$ there is a smooth proper genus-two curve over a
field finitely generated over $\mathbb F_2$, of geometric $2$-rank $f$,
whose Raynaud theta divisor has no nonzero point of any odd prime order.
Consequently the assertion $T(C)$ used in A03 is false on this larger
class of ground fields, even after imposing positive $2$-rank.

\begin{proof}
Choose an irreducible component $Z$ of the nonempty stratum
$\mathcal M_2^f$ and a curve $C_\Omega$ representing its geometric
generic point. Ozman--Pries, Theorem~1.1(2), equivalently
Theorem~8.3, applies with $p=2$, $g=2$, and either $f=1$ or $f=2$:
for every odd prime $\ell$, its theta divisor contains no point of
order $\ell$ \cite[pp.~456,473]{OzmanPries2019}.
These assertions apply to the same geometric generic curve for all
$\ell$; they do not require choosing a different curve for each prime.
The curve and a projective embedding have finitely many coefficients,
so they descend to a finitely generated subfield $K\subset\Omega$.
Smoothness and geometric connectedness descend under this field
extension. The $2$-rank and the absence of geometric theta torsion
are preserved by algebraically closed extension, so this curve over
$K$ has the claimed properties. Any required rational base point can
also be included among the finitely many coefficients.
\end{proof}

This conclusion does not give a counterexample over
$\overline{\mathbb F}_2$: the generic curve has transcendental moduli.
Finding A08 gives the separate finite-field consequence with its
correct quantifiers, namely avoidance of each prescribed finite set
of primes by a curve that may depend on that set.
For the ordinary generic curve, the stronger statement that all
prime-to-$2$ cyclic etale covers are ordinary is the Nakajima result
recorded as Proposition~4 in Yang's Theorem~1.1
\cite[Theorem~1.1]{Yang2021}. No restriction $p>2$ occurs there.
For a generic positive-rank stratum, Yang's Section~3.6.2 and
Corollary~3.6 explicitly retain prime degree; they do not supply a
prime-power version for rank one.

\paragraph{Source qualification.}
The original Nakajima paper is \emph{On generalized Hasse--Witt
invariants of an algebraic curve}, Adv. Stud. Pure Math.
\textbf{2} (1983), 69--88. Its original Proposition~4/Theorem~2 was
not recovered in this follow-up. The generic assertion above is
verified in Yang's paper and, for prime degree, directly in
Ozman--Pries. The latter's Theorem~8.2(3) is not used: its summary of
Pop--Saidi omits their finite-field hypothesis and describes prime
degrees, whereas the cited argument establishes pairwise coprime
orders, not prescribed prime orders. The finite-primary projection
argument in Pop--Saidi, Section~2, is the evidence used below.
\cite[Propositions~2.1,2.3 and their proofs, pp.~112--114]{PopSaidi}
In particular, Yang's stated all-cyclic-degree result also rules
out the sufficient prime-power incidence condition of A11 on a
generic ordinary curve. This stronger consequence is attributed to
the verified statement in Yang, not to an inspection of Nakajima's
unavailable original text. No corresponding all-prime-power claim
is made for the generic rank-one stratum.
There is also a direct modern theorem to cite: Yang's
Theorem~3.5 allows any integer $m$ prime to $p$ and says that a
cyclic cover of a geometric generic pointed curve is new-ordinary
exactly when the associated ramification divisors are Frobenius
stable. In the unpointed etale case every such divisor is zero,
which satisfies Definition~3.3 and Lemma~3.2 automatically.
This verifies the all-odd-degree conclusion at $p=2$ without
claiming access to Nakajima's original paper
\cite[Definition~3.3, Theorem~3.5]{Yang2021}.

